\section{Conclusions and further work}

In retrospect, we managed to define two
self-interpreting type systems, $\SSTLC$ and $\twoLTT$,
meaning that they have a notion of their own ``internal
models'' which they then can use to reinterpret their
own terms inside such models. In doing so, such type
systems get new features which previously were defined
as separate builtin quirks of a system. Basic
metatheoretical properties are proven; for future work
we have the following plans:

\begin{itemize}
    \item We might try making locking discipline in
        Section~\ref{SSTLC} more flexible, allowing
        more than one lock in a context. Most likely,
        this would require us to incorporate modalities
        in the definition of an ``internal model'' type
        and bring us closer to the truly reflective
        system.
    \item While our technique to reinterpret conversion
        rule in Section~\ref{twoLTT} works, a current
        approach to equality witness extraction is a
        \textit{hack}. A proper solution would require us
        to specify witness extraction as a redex and,
        possibly, incorporate certain \textit{algebras
        of reduction strategies} to make this process
        more controllable in proof-relevant settings.
    \item A Categories-with-Families approach to
        specifying internal models here is, in a sense,
        a ``deep'' interface: while all of its
        facilities are properly used in interpretation,
        we have a low-hanging fruit of being able to
        interpret terms in arbitrary CwF context. This
        hints at a direct extension of our modality of
        reinterpretable terms into full-blown contextual
        types \cite{Nanevski:2008:CMT}.
    \item Having implemented contextual types, it then
        becomes possible to handle ``partial''
        reinterpretation of terms: if a term under
        interpretation doesn't use a certain feature of
        a type system, then you shouldn't have to
        specify it in a model! This would bring our
        language closer to Conal Elliott's approach of
        ``Compiling to Categories''
        \cite{Elliott:Compcat2017} and closer to
        an \textit{actual} type system for synthetic
        mathematics.
    \item It would be nice to make equality in $\twoLTT$
        more extensional, most probably by making the type
        system closer to Homotopy type theory
        \cite{hottbook}.
    \item At some moment we would have to pay due
        diligence: provide mechanized proofs of
        metatheoretical properties of systems and an
        implementation of $\twoLTT$ in a practical proof
        assistant.
\end{itemize}
